% chapters/state_of_art/correction_systems.tex
% Sección 2.1: Sistemas de corrección de exámenes (~1.5-2 páginas)
% ============================================================================
%
% GUÍA: Revisar soluciones existentes para corrección de exámenes.
%       Para cada sistema, describir brevemente:
%       - Qué hace y cómo funciona
%       - Ventajas y limitaciones
%       - Comparación con lo que propone el SCDEE
%
%   Sistemas a considerar (investigar y añadir citas ~\cite{...}):
%
%   1. Corrección manual tradicional:
%      - Flujo en papel, sin digitalización
%      - Problemas de escalabilidad, trazabilidad, coordinación
%
%   2. Plataformas de evaluación online (Moodle, Blackboard, Canvas):
%      - Enfocadas a tests automáticos, no exámenes manuscritos
%      - Limitaciones con respuestas de desarrollo libre
%
%   3. Gradescope (Turnitin):
%      - Corrección asistida de exámenes escaneados
%      - OCR para agrupación de respuestas
%      - Referencia obligatoria, es el competidor más directo
%      - Limitaciones: SaaS cerrado, sin autoalojamiento, coste, RGPD
%
%   4. Crowdmark:
%      - Corrección distribuida de exámenes escaneados
%      - Similar a Gradescope pero con enfoque colaborativo
%
%   5. Otros sistemas académicos o investigaciones relevantes:
%      - Buscar en IEEE Xplore, ACM DL artículos sobre
%        "automated exam grading", "handwritten exam processing"
%
%   Cerrar con una tabla comparativa (usar entorno \begin{table}):
%
%   \begin{table}{tab:comparison}{Comparativa de sistemas de corrección}
%     \begin{tabular}{l c c c c c}
%       ...
%     \end{tabular}
%   \end{table}
%
%   Justificar por qué se necesita un sistema propio.
% ============================================================================

% TODO: Investigar, citar y redactar la revisión de sistemas de corrección.
